\documentclass[10pt,a4paper]{amsart}
\usepackage[margin=19mm]{geometry}
\usepackage{amsmath,amssymb,mathtools}
\usepackage[scr=rsfs,cal=euler]{mathalfa}
\usepackage{xfrac}
\usepackage[unicode,hidelinks,bookmarks=false]{hyperref}
\newcommand{\ie}{i.e.\ }
\newcommand{\cf}{cf.\ }
\newcommand{\resp}{respectively\ }
\newcommand{\Subsection}{\subsection}
\providecommand{\nobreakdash}{\nobreak\hbox{-}}
\setlength{\emergencystretch}{2em}
\multlinegap0pt
\usepackage{siunitx}
\usepackage{siunitx}
\usepackage{siunitx}
\usepackage{siunitx}
\usepackage{siunitx}
\usepackage{tikz}
\usepackage{bm}
\newtheorem{lemma}{Lemma}
\newcommand{\param}{\mathit}
\title{Integrating Column Generation and Large Neighborhood Search for Bus Driver Scheduling with Complex Break Constraints}
\author{Lucas Kletzander, Tommaso Mannelli Mazzoli, Nysret Musliu, Pascal Van Hentenryck}
\date{Source: arXiv:2505.02485v1 (CC BY 4.0)}
\begin{document}
\maketitle

\noindent\emph{Source and licence.} Integrating Column Generation and Large Neighborhood Search for Bus Driver Scheduling with Complex Break Constraints. Version arXiv:2505.02485v1 (CC BY 4.0), distributed by arXiv under the Creative Commons Attribution 4.0 International licence, http://creativecommons.org/licenses/by/4.0/, as recorded in the arXiv licence metadata for this exact version and shown on the abstract page https://arxiv.org/abs/2505.02485v1. Full licence notice: \texttt{licenses/arxiv-2505.02485-CC-BY-4.0.txt}.\par\smallskip

\noindent\textit{Editorial scope.} Counted unit: the article's lemma lemma:total\_order (source0.tex lines 202--209). The article's complete original proof of it is delivered with it (source0.tex lines 210--232, one \texttt{proof} environment of 23 lines). No mathematics is written, repaired or re-derived: the statement and the proof are the article's own lines, in the article's own notation, and no retained line is substituted or re-flowed.  No further source-local context is needed: every cross-reference of every retained line resolves inside the two delivered spans.  Nothing was omitted from any retained span, so the package carries no omission record.  No retained line contains a citation, so the wrapper carries no bibliography and no bibliography entry.  No retained line contains a figure environment, an image-inclusion command or drawing code, so no image is delivered and the package declares no asset dependency.  Editorial formatting only, in addition: an AMS \texttt{amsart} preamble that restates the article's own declarations and macros used by the retained lines, namely the environments \texttt{lemma} on separate counters, together with the standard packages those lines need.  The retained environments print in this document as Lemma 1 in order of appearance, whereas the article prints the same items with the numbering of its own full document.  The complete native source of the article as distributed in the arXiv e-print for this exact version is bundled unchanged as \texttt{sources/177588/source0.tex}.
\begin{lemma}\label{lemma:total_order}
Let $\preceq$ be the relation on $L$ defined as follows: for any $\ell_1,\ell_2 \in L$, write $\ell_1\preceq \ell_2$ if either 
\begin{itemize}
    \item $\param{start}_{\ell_1} <\param{start}_{\ell_2} $, or 
    \item $\param{start}_{\ell_1} = \param{start}_{\ell_2}$ and $ \param{tour}_{\ell_1}  \le \param{tour}_{\ell_2}$
\end{itemize}
Then  $\preceq$ is an order relation, and  $(L, \preceq)$ is a totally ordered set.
\end{lemma}

\begin{proof}We first note that in the dataset, each pair $(\param{start}_\ell,\param{tour}_\ell)$ uniquely identifies a bus leg. That is, if two legs share the same start time and tour index, they are the same element of $L$.
\begin{itemize}
        \item \textit{Reflexivity}: For each $\ell \in L$ we have $\param{start}_{\ell}=\param{start}_{\ell}$ and $\param{tour}_{\ell}  = \param{tour}_{\ell}$. So $\ell \preceq \ell $.
   \item  \textit{Antisymmetry} Let $\ell_1,\ell_2 \in L$ be such that $\ell_1\preceq \ell_2$ and $\ell_2\preceq \ell_1$. This implies that $\param{start}_{\ell_1}=\param{start}_{\ell_2}$, $\param{tour}_{\ell_1}  \le \param{tour}_{\ell_2}$, and $\param{tour}_{\ell_2}  \le \param{tour}_{\ell_1}$. So $\ell_1 = \ell_2$.
   \item \textit{Transitivity} Let $\ell_1,\ell_2, \ell_3 \in L$ be such that  $\ell_1 \preceq \ell_2$, and $\ell_2 \preceq \ell_3$. Then there are four cases:
   \begin{enumerate}
       \item  $\param{start}_{\ell_1} <\param{start}_{\ell_2}<\param{start}_{\ell_3}$. Then $\param{start}_{\ell_1} < \param{start}_{\ell_3}$ and therefore $\ell_1 \preceq \ell_3$.
       \item  $\param{start}_{\ell_1} =\param{start}_{\ell_2}<\param{start}_{\ell_3}$ so $\ell_1 \preceq \ell_3$.
              \item  $\param{start}_{\ell_1} <\param{start}_{\ell_2}=\param{start}_{\ell_3}$ so $\ell_1 \preceq \ell_3$.
       
       \item $\param{start}_{\ell_1} =\param{start}_{\ell_2}=\param{start}_{\ell_3}$ and $\param{tour}_{\ell_1}\le \param{tour}_{\ell_2} \le \param{tour}_{\ell_3}$. This implies $\ell_1 \preceq \ell_3$.
   \end{enumerate}
   \item \textit{Totally ordered}: For any $\ell_1, \ell_2 \in L$, exactly one of the following holds:  
   \begin{itemize}
       \item   $\param{start}_{\ell_1} < \param{start}_{\ell_2}$,  
   \item  $\param{start}_{\ell_1} > \param{start}_{\ell_2}$, or  
   \item $\param{start}_{\ell_1} = \param{start}_{\ell_2}$
   \end{itemize}
   In the first two cases, $\ell_1 \preceq \ell_2$ or $\ell_2 \preceq \ell_1$ follows directly.  If $\param{start}_{\ell_1} = \param{start}_{\ell_2}$, then $\param{tour}_{\ell_1} \leq \param{tour}_{\ell_2}$ or $\param{tour}_{\ell_1} \geq \param{tour}_{\ell_2}$. Thus, $\ell_1 \preceq \ell_2$ or $\ell_2 \preceq \ell_1$.  
   Hence, $(L, \preceq)$ is totally ordered.
 \qedhere
    \end{itemize}
\end{proof}

\end{document}
